\documentclass[10pt,a4paper]{amsart}
\usepackage[margin=19mm]{geometry}
\usepackage{amsmath,amssymb,mathtools}
\usepackage[scr=rsfs,cal=euler]{mathalfa}
\usepackage{xfrac}
\usepackage[unicode,hidelinks,bookmarks=false]{hyperref}
\newcommand{\ie}{i.e.\ }
\newcommand{\cf}{cf.\ }
\newcommand{\resp}{respectively\ }
\newcommand{\Subsection}{\subsection}
\providecommand{\nobreakdash}{\nobreak\hbox{-}}
\setlength{\emergencystretch}{2em}
\multlinegap0pt
\usepackage{amssymb}
\usepackage[shortlabels]{enumitem}
\newtheorem{lem}{Lemma}
\def\A{\mathcal A}
\def\B{\mathcal B}
\def\C{\mathcal C}
\def\DD{\Delta}
\def\ss{\subset}
\title{Moving rectangular sofas in planar and spatial corridors}
\author{Oleg Mushkarov, Nikolai Nikolov}
\date{Source: arXiv:2604.01174v1 (CC BY 4.0)}
\begin{document}
\maketitle

\noindent\emph{Source and licence.} Moving rectangular sofas in planar and spatial corridors. Version arXiv:2604.01174v1 (CC BY 4.0), distributed by arXiv under the Creative Commons Attribution 4.0 International licence, http://creativecommons.org/licenses/by/4.0/, as recorded in the arXiv licence metadata for this exact version and shown on the abstract page https://arxiv.org/abs/2604.01174v1. Full licence notice: \texttt{licenses/arxiv-2604.01174-CC-BY-4.0.txt}.\par\smallskip

\noindent\textit{Editorial scope.} Counted unit: the article's lem l3 (source0.tex lines 405--409). The article's complete original proof of it is delivered with it (source0.tex lines 411--457, one \texttt{proof} environment of 47 lines). No mathematics is written, repaired or re-derived: the statement and the proof are the article's own lines, in the article's own notation, and no retained line is substituted or re-flowed.  Retained uncounted as the source-local context the proof needs: lines 362--365.  Nothing was omitted from any retained span, so the package carries no omission record.  No retained line contains a citation, so the wrapper carries no bibliography and no bibliography entry.  No retained line contains a figure environment, an image-inclusion command or drawing code, so no image is delivered and the package declares no asset dependency.  Editorial formatting only, in addition: an AMS \texttt{amsart} preamble that restates the article's own declarations and macros used by the retained lines, namely the environments \texttt{lem} on separate counters, together with the standard packages those lines need.  The retained environments print in this document as Lemma 1 in order of appearance, whereas the article prints the same items with the numbering of its own full document.  The complete native source of the article as distributed in the arXiv e-print for this exact version is bundled unchanged as \texttt{sources/197641/source0.tex}.
\begin{lem}\label{l2} A rectangle $R\ss \A_1$ with side lengths
$c$ and $d$ (not necessarily $c\ge d$) can be moved into $\B_3$ by
a rotation around the corner $C$ if and only if $d\le m(a,b,c).$
\end{lem}

\begin{lem}\label{l3} A rectangle $R\ss \A_1$ with side lengths
$c$ and $d,$ $c\ge a\vee b,$ can be moved into $\B_3$ by an
anti-rotation around the corner $C$ if and only if $c\le
\widetilde{m}(a,b,d).$
\end{lem}

\begin{proof}
We first prove the \emph{only if part} of the lemma. Let a
rectangle $R\ss \A_1$ with side lengths $c$ and $d,$ $c\ge a\vee
b,$ can be moved into $\B_3$ by an anti-rotation around the corner
$C$. Note that for any $t\in \DD$ there is a moment when the angle
between a side of $R$ with length $d$ and $Ox$ is $t$. In
particular, at some moment such a side of $R$ is perpendicular to
$OC$ and therefore $c^2\le a^2+b^2$ which implies that $t_1\le
t_2$. We note that for $t\in [0,t_1]$ the rectangle $R$ lies in
the corridor $\A_1$, whereas for $t\in [t_2,\pi/2]$ it lies in the
corridor $\B_3$. Hence it is enough to consider only the positions
of $R$ where $t\in \widetilde{\DD}$. For any such $t$ there is a
moment when the sides of $R$ with length $d$ are parallel to the
tangent to the circle $k(C;c)$ with slope $-\tan t$ and we
conclude as in the proof of the \emph{only if part} of Lemma
\ref{l2} that $c\le \widetilde{m}(a,b,d)$.

To prove the \emph{if part } of the lemma suppose that $c\le
\widetilde{m}(a,b,d)$ and note that this inequality can be written
as
\begin{equation}\label{2}
d\le g_{(a,b,c)}(t)=\frac{a\sin t+b\cos t-c}{\sin t\cos t} .
\end{equation}
It is easy to check that $g_{(a,b,d)}(t)$ is the length of the
segment which the tangent to the circle $k(C;d)$ and slope $-\tan
t$ cuts from the first quadrant. In particular, we have from
(\ref{2}) that
$$ d\le g_{(a,b,c)}(t_1)=\frac{c(a-\sqrt{c^2-b^2)}}{b}\leq a $$
and
$$ d\le g_{(a,b,c)}(t_2)=\frac{c(b-\sqrt{c^2-a^2)}}{a}\leq b $$
since $c\ge a\vee b .$

Let now $R=KLMN$ be a rectangle with side lengths $c$ and $d,$
$c\ge a\vee b,$ which is in a good position with $M=O(0,0)$ and
$\overrightarrow{ML} = c\overrightarrow{e_2} $. For any $t\in
\widetilde{\DD}$ consider the rectangle $R_t=K_tL_tM_tN_t$ with
vertices $K_t=(d\cos t+c\sin t, c\cos t) , L_t=(c\sin t, c\cos
t+d\sin t) , M_t=(0,d\sin t) , N_t=(d\cos t,0) $. It is easy to
check that the function $x(t)=d\cos t+c\sin t$ is increasing on
the segment $[0,t_1]$ since $d\le a$ and $x(t_1)\le a $ since
$c\le f_{(a,b,d)}(t_1)$. Hence $K_t\in \A_1$ for any $t\in
[0,t_1]$ which implies that $R_t \ss \A$ for $t\in [0,t_1]$.
Similarly $R_t \ss \B$ for $t\in [t_2, \pi/2)$. It follows also
that $R_t \ss \C$ for $t\in [t_1,t_2],$ since the inequality $c\le
f_{(a,b,d)}(t) $ means that the point $C$ is at a distance at
least $d$ from the line $M_tN_t$.
\end{proof}

\end{document}
